% ============================================================
%  English Version — sync from main_zh.tex after Chinese editing
%  DO NOT edit directly; wait for confirmed Chinese draft first
% ============================================================
\documentclass[12pt]{article}

\usepackage{graphicx}
\usepackage{amsmath, amssymb}
\usepackage{booktabs}
\usepackage{hyperref}
\usepackage{geometry}
\usepackage{cite}
\geometry{a4paper, margin=2.5cm}

% ---------- Paper Info ----------
\title{\textbf{DiffAug: Conditional Diffusion with Physical-Consistent Projection\\
for Traffic Speed Prediction under Sparse Supervision}}
\author{
  Yijin Yan$^{1}$\quad Guojian Zou$^{1}$\\[4pt]
  $^{1}$ [Institution]\\
  \texttt{[email]}
}
\date{March 2026}

\begin{document}

\maketitle

% ============================================================
\begin{abstract}
% TODO: 200--250 words (sync from Chinese draft)
Traffic speed prediction is a critical component of intelligent transportation systems (ITS).
In real-world deployments, two fundamental challenges arise:
\textbf{data scarcity}—private road networks accumulate far fewer historical samples than public
benchmarks—and \textbf{holiday distribution shift}, where traffic patterns during holidays
deviate systematically from weekday distributions.

We propose \textbf{DiffAug}, a generative multi-view augmentation framework built on a
conditional denoising diffusion probabilistic model (DDPM).
Given $K$ augmented speed views sampled from the conditional distribution of
historical $(v, q, \rho)$ triplets,
we apply a \textit{physical-consistency projection}
$\tilde{\rho} = q/(\tilde{v}+\varepsilon)$
to strictly enforce the traffic flow conservation equation $q = \tilde{\rho}\cdot\tilde{v}$.
A Beta-interpolation TSMixup step further covers the continuum between original and generated views,
and real observations act as a soft anchor to stabilize supervised training.

Experiments on the private MSST-CL dataset ($N=30$ nodes, 15-min granularity)
demonstrate that DiffAug improves MAE/RMSE over no-augmentation baselines
while maintaining near-perfect distributional fidelity
(KL $<0.01$, WD $<0.5$\,km/h) between generated and real speed distributions.

\textbf{Keywords:} Traffic speed prediction, diffusion model, data augmentation,
physical consistency, sparse supervision
\end{abstract}

% ============================================================
\section{Introduction}
\label{sec:intro}

% --- Background and Motivation ---
Traffic speed prediction is a foundational task in intelligent transportation systems (ITS),
underpinning route planning, signal control, and congestion alerting.
Spatio-temporal deep learning models, including graph neural networks (GNNs) and
Transformers, have achieved remarkable results on public
benchmarks~\cite{stgformer,staeformer}.

However, real-world private road-network deployments face two typical difficulties:
\begin{enumerate}
  \item \textbf{Data scarcity.}
    Private road networks accumulate limited historical data—often fewer than 13,000
    training windows, roughly one-tenth of public datasets such as PEMS04/08—making
    models prone to overfitting.
  \item \textbf{Holiday distribution shift.}
    Holiday traffic volumes surge and time-of-day patterns deviate systematically
    from weekday norms, leading to poor generalization on holiday periods.
\end{enumerate}

Existing data augmentation approaches broadly fall into three categories:
(1) \textit{structural augmentation} (masking, dropout, temporal cropping), which lacks
physical-semantic constraints;
(2) \textit{interpolation augmentation} (Mixup~\cite{mixup}), which is simple but cannot
guarantee traffic-physical validity;
(3) \textit{generative augmentation} (STGAN~\cite{stgan}), which is difficult to control
in terms of distributional alignment.

Our key insight is that \textit{effective traffic augmentation views must simultaneously satisfy:
(a) distributional fidelity—closely matching real traffic speed statistics,
(b) sample diversity—introducing speed scenarios unseen in the weekday training set, and
(c) physical legality—satisfying the flow conservation constraint
$q = \tilde{\rho}\cdot\tilde{v}$.}

We therefore propose \textbf{DiffAug}, with the following contributions:
\begin{itemize}
  \item We position a conditional DDPM as a \textbf{data generator} decoupled from the
    prediction backbone, allowing independent optimization of generation quality.
  \item We design a \textbf{physical-consistency projection} module that enforces
    speed–flow–density conservation after diffusion sampling.
  \item We introduce a \textbf{distribution-quality filtering} strategy that uses
    KL divergence and Wasserstein distance to retain only high-fidelity augmented views.overleaf_session2=s%3Ag1zUXB1GwD9EbKRlgZsmNm2i_8IZws9B.%2BbhIqq74wmBv5Uw%2BddYuRn1lxYcg2sqthbfG3mN6RW0; GCLB=CJqY4qTJycTBBhAD
\section{Related Work}
\label{sec:related}

\subsection{Spatio-Temporal Traffic Prediction}overleaf_session2=s%3Ag1zUXB1GwD9EbKRlgZsmNm2i_8IZws9B.%2BbhIqq74wmBv5Uw%2BddYuRn1lxYcg2sqthbfG3mN6RW0; GCLB=CJqY4qTJycTBBhAD
% TODO: STGformer, STAEformer, STGCN, DCRNN, etc.
% Reference: doc/need/相关工作检索与对比模型清单.md

\subsection{Data Augmentation for Traffic Prediction}
% TODO: Mixup, CutMix, STGAN, STAug
% Emphasize: physical constraint distinguishes our work

\subsection{Diffusion Models for Time-Series Generation}
% TODO: DDPM, DDIM, DiffSTG, TimeDiff
% Emphasize: "generator, not predictor" positioning

% ============================================================
\section{Problem Formulation}
\label{sec:problem}

Consider a road network with $N$ nodes observed over $T$ time steps at 15-minute
granularity. The traffic state at each step is described by speed $v\in\mathbb{R}^N$,
flow $q\in\mathbb{R}^N$, and density $\rho\in\mathbb{R}^N$, satisfying the
fundamental traffic flow equation:
\begin{equation}
  q = \rho \cdot v.
  \label{eq:traffic_fundamental}
\end{equation}

Given a historical window $\mathbf{X} = \{(v_t, q_t, \rho_t)\}_{t=1}^{L}$ ($L=12$),
the goal is to predict the future $H$-step speed
$\hat{\mathbf{v}} = \{\hat{v}_t\}_{t=L+1}^{L+H}$ ($H=12$).

Due to limited training data ($\sim$12,707 windows),
we seek to construct $K$ physically valid augmented views
$\{\tilde{v}^{(k)}\}_{k=1}^K$ that share the same future label $\mathbf{Y}$,
effectively expanding the training set.

% ============================================================
\section{Methodology}
\label{sec:method}

\subsection{Overall Framework}
% TODO: Insert architecture figure

DiffAug consists of three modules:
\begin{enumerate}
  \item \textbf{Conditional DDPM Generator} (Sec.~\ref{sec:ddpm}):
    samples $K$ augmented speed views conditioned on historical traffic state;
  \item \textbf{Physical-Consistency Projection} (Sec.~\ref{sec:projection}):
    projects generated speeds to the feasible state space satisfying the
    conservation equation;
  \item \textbf{TSMixup Interpolation \& Weighted Training} (Sec.~\ref{sec:mixup}):
    Beta interpolation covers intermediate augmentation states;
    a real-view soft anchor stabilizes supervised training.
\end{enumerate}

\subsection{Conditional Diffusion Generator}
\label{sec:ddpm}

We adopt the DDPM~\cite{ddpm} forward noising process:
\begin{equation}
  q(\mathbf{x}_s | \mathbf{x}_0) =
  \mathcal{N}\!\left(\mathbf{x}_s;\, \sqrt{\bar{\alpha}_s}\,\mathbf{x}_0,\,
  (1-\bar{\alpha}_s)\mathbf{I}\right),
\end{equation}
where $\mathbf{x}_0 = v_{1:L}$ is the historical speed sequence, and the condition
$c$ encodes historical flow $q_{1:L}$, time-of-day (ToD), and day-of-week (DoW).

The denoising network $\epsilon_\theta(\mathbf{x}_s, s, c)$ uses a Patch Transformer
with ALiBi positional encoding~\cite{alibi}.
Training minimizes the standard denoising MSE:
\begin{equation}
  \mathcal{L}_{\text{denoise}} =
  \mathbb{E}_{\mathbf{x}_0, \epsilon, s}\!\left[
    \|\epsilon - \epsilon_\theta(\mathbf{x}_s, s, c)\|^2
  \right].
\end{equation}

% TODO: Add statistical loss (stat_loss) formulation and motivation

\subsection{Physical-Consistency Projection}
\label{sec:projection}

After sampling augmented speed $\tilde{v}^{(k)}$, the corresponding density is
computed to enforce conservation:
\begin{equation}
  \tilde{\rho}^{(k)} = \frac{q}{\tilde{v}^{(k)} + \varepsilon},
  \label{eq:projection}
\end{equation}
where $\varepsilon$ is a numerical stability term and $q$ is the observed flow
(treated as fixed).
This guarantees that $(q, \tilde{v}^{(k)}, \tilde{\rho}^{(k)})$ strictly satisfies
Eq.~\eqref{eq:traffic_fundamental}.

\subsection{TSMixup Interpolation and Weighted Training}
\label{sec:mixup}

To cover the continuum from the original to the generated view, we apply
Beta interpolation to each augmented view:
\begin{equation}
  v_{\text{aug}}^{(k)} = \lambda^{(k)}\, v + (1-\lambda^{(k)})\, \tilde{v}^{(k)},
  \quad \lambda^{(k)} \sim \text{Beta}(\alpha, \alpha).
  \label{eq:mixup}
\end{equation}

The training loss combines the real-view loss and augmented-view loss:
\begin{equation}
  \mathcal{L} = \mathcal{L}(v, \mathbf{Y}) +
  \frac{\alpha_{\text{aug}}}{K}\sum_{k=1}^K \mathcal{L}(v_{\text{aug}}^{(k)}, \mathbf{Y}),
  \label{eq:loss}
\end{equation}
where $\alpha_{\text{aug}} < 1$ ensures real observations remain the primary supervision
signal.

% ============================================================
\section{Experiments}
\label{sec:exp}

\subsection{Dataset and Setup}

\paragraph{MSST-CL Dataset.}
% TODO: dataset description
Private road network; $N=30$ nodes; 15-minute granularity;
train/val/test split 60\%/20\%/20\%;
$\sim$12,707 training windows ($L=12$, $H=12$).
Contains holiday periods with significant distribution shift.

\paragraph{Baselines.}
% TODO: list baselines
STGformer~\cite{stgformer},
STGCN~\cite{stgcn},
DCRNN~\cite{dcrnn},
STAEformer~\cite{staeformer},
DiffSTG~\cite{diffstg}.

\paragraph{Metrics.}
MAE, RMSE, MAPE (averaged over $H=12$ future steps).

\subsection{Main Results}
% TODO: Insert main results table

\subsection{Ablation Study}

Table~\ref{tab:ablation} shows the contribution of each component.
\begin{table}[ht]
  \centering
  \caption{Ablation study on the MSST-CL dataset (MAE / RMSE).}
  \label{tab:ablation}
  \begin{tabular}{lcccc}
    \toprule
    \textbf{Variant} & \textbf{Strategy} & \textbf{KL}$\downarrow$
      & \textbf{MAE}$\downarrow$ & \textbf{RMSE}$\downarrow$ \\
    \midrule
    No augmentation  & —           & —      & 4.4045 & 6.6519 \\
    e4               & stat\_loss  & 0.0022 & 4.4225 & 6.6494 \\
    e1               & adaLN       & —      & 4.4473 & 6.6907 \\
    Baseline aug.    & —           & —      & 4.4413 & 6.6859 \\
    \midrule
    \textbf{DiffAug (ours)} & e3e4 + filtering
      & \textbf{TODO} & \textbf{TODO} & \textbf{TODO} \\
    \bottomrule
  \end{tabular}
\end{table}

\subsection{Generation Quality Analysis}
% TODO: Insert KL, WD, intra_std analysis figures/tables

% ============================================================
\section{Conclusion}
\label{sec:conclusion}

We presented DiffAug, a conditional diffusion-based generative multi-view augmentation
framework for traffic speed prediction on sparse private road networks.
By combining physical-consistency projection with TSMixup interpolation,
DiffAug produces augmented views that are both distributional-fidelity-compliant
and sufficiently diverse to expose the prediction backbone to previously unseen speed
scenarios.
Experiments on the private MSST-CL dataset confirm that DiffAug effectively mitigates
overfitting and improves generalization under holiday distribution shift.

\paragraph{Limitations and Future Work.}
The current filtering strategy relies on offline precomputation; future work will
explore online adaptive filtering.
End-to-end joint training of the diffusion generator and the prediction backbone
also warrants further investigation.

% ============================================================
\bibliographystyle{plain}
\bibliography{refs}

% ============================================================
\appendix
\section{Appendix: Diffusion Variant Hyperparameters}
% TODO: Hyperparameter table for all variants

\end{document}
